\scanpage{48}
\uline{Modèles de Néron et groupes p-divisibles.}

Dans ce qui suit $S$ est un trait de point générique $\handwritten{\eta}$ et de point
fermé $s$. Un $S$-groupe p-divisible $(G_n,\ n\handwritten{\in}N)$ sera aussi désigné par
le symbole $\underset{\handwritten{\scriptscriptstyle\sim\!\sim}}{G}$ . On \deleted{sé}\inserted{\handwritten{dés}}igne par $p$ un nombre premier fixé.

\uline{Théorème} : Soient $G$ un $S$-schéma en groupes, $S$-néronien,\handwritten{[}dont la fibre générique est extension d'une variété abélienne par un tore\handwritten{]} et $\underset{\handwritten{\scriptscriptstyle\sim\!\sim}}{H}$ un $S$-groupe
% slanted
\margin{\handwritten{[ ] inutile}}
p-divisible. Alors tout morphisme $u_{\handwritten{\eta}}:(\underset{\handwritten{\scriptscriptstyle\sim\!\sim}}{H})_{\handwritten{\eta}}\handwritten{\to} G_{\handwritten{\eta}}$ se prolonge de manière unique en un $S$-morphisme $u:\underset{\handwritten{\scriptscriptstyle\sim\!\sim}}{H}\handwritten{\to} G$ .

\uline{Lemme} 1: Supposons $S$ hensélien. Soient $G$ un $S$-schéma en groupes , $S$-néronien \deleted{et}\handwritten{,} $M$ un $S$-schéma en groupes fini et plat , $\overset{\handwritten{\sim}}{M}$ sa partie connexe,
\deleted{$u_K$} $u_{\handwritten{\eta}}:M_{\handwritten{\eta}}\handwritten{\to} G_{\handwritten{\eta}}$ un homomorphisme. Alors , pour que $u_{\handwritten{\eta}}$ se prolonge en
un $S$-morphisme $u:M\handwritten{\to} G$, il faut et il suffit que $u_{\handwritten{\eta}}\handwritten{\mid}\overset{\handwritten{\sim}}{M}_{\handwritten{\eta}}$ se prolonge
en un $S$-morphisme $\overset{\handwritten{\sim}}{u}:\overset{\handwritten{\sim}}{M}\handwritten{\to} G$ .

% slanted
\margin{\handwritten{$G$ commutatif?}}
La nécessité est clai\deleted{\unsure{t}}\unsure{r|t}e. Prouvons la suffisance. Considérons le groupe
produit $M\times_S G$ et le morphisme
\[
\begin{aligned}
v_{\handwritten{\eta}}:M_{\handwritten{\eta}}\handwritten{\times}_{\handwritten{\eta}}G_{\handwritten{\eta}}&\handwritten{\to} G_{\handwritten{\eta}}\\
(m,g)&\handwritten{\mapsto} u_{\handwritten{\eta}}(m)+g
\end{aligned}
\]
Il nous suffit de montrer que $v_{\handwritten{\eta}}$ se prolonge en un morphisme $v:M\handwritten{\times}_S G\handwritten{\to} G$,
car la restriction de $v$ au premier facteur $M$ fournira le prolongement
de $u$ . Soit $P_{\handwritten{\eta}}$ le noyau de $v_{\handwritten{\eta}}$ et $P$ l'adhérence schématique de $P_{\handwritten{\eta}}$ dans
$M\handwritten{\times}_S G$. Le lemme 1 résulte de la propriété Néronienne de $G$ et du lemme suivant:

\uline{Lemme 2}: Le $S$-schéma en groupes $M\handwritten{\times}_S G/P$ \inserted{\handwritten{existe et}} est lisse sur $S$.

Vu l'hypothèse faite sur $u_{\handwritten{\eta}}$, le morphisme $\overset{\handwritten{\sim}}{v}_{\handwritten{\eta}}=v_{\handwritten{\eta}}\handwritten{\mid}\overset{\handwritten{\sim}}{M}_{\handwritten{\eta}}\handwritten{\times}_{\handwritten{\eta}}G_{\handwritten{\eta}}$ se prolonge en un $S$-morphisme
\[
\begin{aligned}
\overset{\handwritten{\sim}}{v}:\overset{\handwritten{\sim}}{M}\handwritten{\times}_S G&\handwritten{\to} G\\
(\overset{\handwritten{\sim}}{m},g)&\handwritten{\mapsto} \overset{\handwritten{\sim}}{u}(m)+g
\end{aligned}
\]
de sorte que l'on a la suite exacte :
\[
0\handwritten{\to}\operatorname{Ker}\overset{\handwritten{\sim}}{v}\handwritten{\to}\overset{\handwritten{\sim}}{M}\handwritten{\times}_S G\handwritten{\to} G\handwritten{\to}0 \tag*{\makebox[0pt][r]{\makebox[\displaywidth][l]{(1)}}}
\]
\margin{\handwritten{$\checkmark$}}
% archive sheet number: 47
